% Figure: the spectral basis of a selective surface. The 5800 K solar
% spectrum (short wavelength) and the 300 K surface spectrum (long
% wavelength), each normalised to unit peak, with a step-function
% spectral absorptance overlaid: high in the visible, low in the infrared.
% Curves are Planck shapes evaluated at the two temperatures, normalised.
\begin{figure}[htbp]
\centering
\begin{tikzpicture}
\begin{axis}[
  width=12cm, height=7cm,
  xmode=log,
  xlabel={wavelength $\lambda$ ($\mu$m)},
  ylabel={normalised spectral power / absorptance},
  xmin=0.2, xmax=40,
  ymin=0, ymax=1.15,
  axis lines=left,
  tick label style={font=\scriptsize},
  label style={font=\small},
  legend style={font=\scriptsize, at={(0.98,0.98)}, anchor=north east},
  legend cell align=left,
  legend entries={sun (5800 K),surface (300 K),selective absorptance},
  clip=false,
]
% Normalised Planck-shape proxies: use lognormal bumps peaking at Wien peaks
% solar peak ~ 0.5 um, surface peak ~ 9.7 um. Use exp(-((ln(l/lp))^2)/(2 s^2)).
\addplot[engamber, very thick, domain=0.2:40, samples=120]
  {exp(-((ln(x/0.5))^2)/(2*0.32))};
\addplot[engblue, very thick, domain=0.2:40, samples=120]
  {exp(-((ln(x/9.7))^2)/(2*0.30))};
% step-function selective absorptance: high below ~2.5 um, low above
\addplot[engred, thick, dashed] coordinates {(0.2,0.95) (2.5,0.95) (2.5,0.10) (40,0.10)};
\draw[enggray, dotted] (axis cs:2.5,0) -- (axis cs:2.5,1.15);
\node[font=\scriptsize, enggray, anchor=south, rotate=90] at (axis cs:2.5,0.30) {cut-off};
\end{axis}
\end{tikzpicture}
\caption[The spectral basis of a selective surface]{Why a surface can
absorb sunlight and yet emit weakly, or the reverse. The incoming solar
spectrum peaks in the visible near $0.5\,\mu$m, while a surface near room
temperature emits in the infrared near $10\,\mu$m; the two bands barely
overlap. A spectrally selective coating exploits the gap by carrying a
high absorptance below a cut-off wavelength and a low one above it (dashed
step). A solar absorber wants the profile drawn here, high in the solar
band and low in its own emission band, so it takes in sunlight and holds
the heat; a surface that must stay cool in sun wants the mirror image, low
solar absorptance and high infrared emittance. The single design number is
the ratio $\alpha_s/\epsilon$ of the band-averaged values, which sets the
equilibrium temperature in sunlight.}
\label{fig:vol04ch05:selective-surface}
\end{figure}
